\documentclass[10pt]{book}
\usepackage[sectionbib]{natbib}
\usepackage{array,epsfig,fancyhdr,rotating}
\usepackage{hyperref}
\usepackage{xr-hyper}
\externaldocument{profile_sieve_geometric_paper}
\usepackage{amsmath}
\usepackage{amssymb}
\usepackage{amsfonts}
\usepackage{amsthm}
\usepackage{bm}
\usepackage{graphicx}
\graphicspath{{writing_samples/figures/}{figures/}}

\textwidth=31.9pc
\textheight=46.5pc
\oddsidemargin=1pc
\evensidemargin=1pc
\headsep=15pt
\headheight=28pt
\topmargin=.6cm
\parindent=1.7pc
\parskip=0pt

\numberwithin{equation}{section}
\setcounter{page}{1}
\newtheorem{theorem}{Theorem}
\newtheorem{lemma}{Lemma}
\newtheorem{corollary}{Corollary}
\newtheorem{proposition}{Proposition}
\newtheorem{assumption}{Assumption}
\theoremstyle{definition}
\newtheorem{definition}{Definition}
\newtheorem{example}{Example}
\newtheorem{remark}{Remark}

\newcommand{\ind}{\mathbf 1}
\newcommand{\R}{\mathbb R}

\pagestyle{fancy}
\def\n{\noindent}
\lhead[\fancyplain{} \leftmark]{}
\chead[]{}
\rhead[]{\fancyplain{}\rightmark}
\cfoot{\thepage}

\begin{document}
\raggedbottom

\renewcommand{\baselinestretch}{2}

\markright{ \hbox{\footnotesize\rm Statistica Sinica: Supplement
%{\footnotesize\bf 24} (201?), 000-000
}\hfill\\[-13pt]
\hbox{\footnotesize\rm
}\hfill }

\markboth{\hfill{\footnotesize\rm SIYAN DONG} \hfill}
{\hfill {\footnotesize\rm MULTIPLE NUISANCE BREAKS AND PREDICTABILITY} \hfill}

\renewcommand{\thefootnote}{}
$\ $\par \fontsize{12}{14pt plus.8pt minus .6pt}\selectfont

\centerline{\large\bf MULTIPLE NUISANCE BREAKS AND}
\vspace{2pt}
\centerline{\large\bf SPURIOUS PREDICTABILITY UNDER}
\vspace{2pt}
\centerline{\large\bf PERSISTENT REGRESSORS}
\vspace{.25cm}
\centerline{Siyan Dong}
\vspace{.4cm}
\centerline{\it KU}
\vspace{.55cm}
\centerline{\bf Supplementary Material}
\vspace{.55cm}
\fontsize{9}{11.5pt plus.8pt minus .6pt}\selectfont
\noindent
This supplement gives the technical details behind the block-sieve empirical likelihood construction in the main manuscript. Section S1 records the finite-sample block projection identities using the Moore--Penrose inverse. Section S2 proves the stable-score decomposition and the conditional omitted-break failure theorem. Section S3 proves the known-partition Wilks theorem under projection-bias negligibility rather than full root-\(T\) undersmoothing. Section S4 gives the high-level fixed-order estimated-score likelihood transfer. Section S5 proves the shifted-array local-alternative calculation. Section S6 records implementation details for the reported simulations and the empirical application.
\par

\setcounter{section}{0}
\setcounter{equation}{0}
\renewcommand{\thesection}{S\arabic{section}}
\renewcommand{\theequation}{\thesection.\arabic{equation}}
\renewcommand{\thelemma}{S\arabic{lemma}}
\renewcommand{\theproposition}{S\arabic{proposition}}
\renewcommand{\thecorollary}{S\arabic{corollary}}
\renewcommand{\thetheorem}{S\arabic{theorem}}
\renewcommand{\theremark}{S\arabic{remark}}

\fontsize{12}{14pt plus.8pt minus .6pt}\selectfont

For a scalar array \(\bm a=(a_1,\ldots,a_T)'\), write
\[
    \|\bm a\|_T^2
    =
    T^{-1}\sum_{t=1}^T a_t^2,
    \qquad
    \|\bm a\|_\infty=\max_{t\leq T}|a_t|.
\]
Unless stated otherwise, \(O_p(\cdot)\) and \(o_p(\cdot)\) for growing
vectors and matrices are in Euclidean and spectral norm, respectively.

\section{Block Projection Algebra}

Fix a partition \(\Lambda=(k_1,\ldots,k_q)\). Let \(\bm Q_{K,\Lambda}\) and \(\bm M_{K,\Lambda}\) be defined as in (\ref*{eq:block-q}) and (\ref*{eq:block-m}). Throughout this supplement, all projection equalities are equalities in the empirical Hilbert space \((\R^T,T^{-1}\bm a'\bm b)\).

\begin{lemma}[Block projection identities]
\label{lem:s-block-identities}
For any partition \(\Lambda\),
\[
    \bm Q_{K,\Lambda}'\bm M_{K,\Lambda}=\bm0,
    \qquad
    \bm M_{K,\Lambda}'=\bm M_{K,\Lambda},
    \qquad
    \bm M_{K,\Lambda}^2=\bm M_{K,\Lambda}.
\]
For every \(\bm c\) conformable with \(\bm Q_{K,\Lambda}\),
\[
    (\bm Q_{K,\Lambda}\bm c)'\bm M_{K,\Lambda}\bm w=0.
\]
\end{lemma}

\begin{proof}
Let \(\bm\Pi_{K,\Lambda}=\bm Q_{K,\Lambda}(\bm Q_{K,\Lambda}'\bm Q_{K,\Lambda})^\dagger\bm Q_{K,\Lambda}'\). This is the orthogonal projector onto the column space of \(\bm Q_{K,\Lambda}\). Therefore \(\bm\Pi_{K,\Lambda}'=\bm\Pi_{K,\Lambda}\), \(\bm\Pi_{K,\Lambda}^2=\bm\Pi_{K,\Lambda}\), and \(\bm Q_{K,\Lambda}'(\bm I_T-\bm\Pi_{K,\Lambda})=\bm0\). Since \(\bm M_{K,\Lambda}=\bm I_T-\bm\Pi_{K,\Lambda}\), the displayed identities follow. The final equality is obtained by multiplying \(\bm Q_{K,\Lambda}'\bm M_{K,\Lambda}=\bm0\) by \(\bm c\) and \(\bm w\).
\end{proof}

\section{Stable-Score Failure}

This section proves Proposition \ref*{prop:stable-decomp} and Theorem \ref*{thm:failure}. Let \(\Lambda=0\) denote the no-break partition and write \(\bm M_{K,0}\) for the corresponding stable-nuisance residual-maker.

\begin{proof}[Proof of Proposition \ref*{prop:stable-decomp}]
Under \(H_0:\beta=\beta_0\),
\[
    \bm Y-\beta_0\bm X_{lag}=\bm\mu+\bm U.
\]
Multiplying by the stable residualized weight gives the exact identity
\[
    T^{-1/2}(\bm Y-\beta_0\bm X_{lag})'\bm M_{K,0}\bm w
    =
    T^{-1/2}\bm U'\bm M_{K,0}\bm w
    +
    T^{-1/2}\bm\mu'\bm M_{K,0}\bm w.
\]
The second term is the omitted-break component. No stable-sieve approximation remainder is needed because \(B_T\) is defined using the exact nuisance vector \(\bm\mu\).
\end{proof}

\begin{proof}[Proof of Theorem \ref*{thm:failure}]
The scalar EL expansion gives
\[
    \ell_T^{\mathrm{st}}(\beta_0)
    =
    \frac{\{S_T^{\mathrm{st}}(\beta_0)\}^2}{V_T^{\mathrm{st}}}
    +o_p(1).
\]
By Proposition \ref*{prop:stable-decomp}, \(S_T^{\mathrm{st}}(\beta_0)=A_T^{\mathrm{st}}+B_T\). With
\[
    a_T=\frac{A_T^{\mathrm{st}}}{(V_T^{\mathrm{st}})^{1/2}},
    \qquad
    b_T=\frac{B_T}{(V_T^{\mathrm{st}})^{1/2}},
\]
the expansion is \(\ell_T^{\mathrm{st}}(\beta_0)=(a_T+b_T)^2+o_p(1)\). If \(a_T=O_p(1)\) and \(|b_T|\to_p\infty\), then \(|a_T+b_T|\geq |b_T|-|a_T|\to_p\infty\), so the statistic diverges in probability. Under joint stable convergence \((a_T,b_T)\xrightarrow{\mathrm{st}}(Z,B)\), the continuous mapping theorem gives \((a_T+b_T)^2\Rightarrow (Z+B)^2\). Since \(Z\) is conditionally standard normal given the limiting sigma-field, \((Z+B)^2\mid B\sim\chi_1^2(B^2)\).
\end{proof}

\begin{remark}[Scope of omitted-break divergence]
The divergence conclusion is conditional on the stable empirical-likelihood quadratic expansion. With a fixed omitted nuisance break, \(B_T\) may be of order \(\sqrt T\), and the null-style scalar expansion need not hold automatically. A primitive fixed-break divergence proof therefore requires a separate expansion under the misspecified stable nuisance projection.
\end{remark}

\section{Known-Partition Wilks Theorem}

Set \(\Lambda=\Lambda_0\) and write
\[
    \bm v_0=\bm M_{K,\Lambda_0}\bm w,
    \qquad
    \widehat{\bm u}_0(\beta_0)
    =\bm M_{K,\Lambda_0}(\bm Y-\beta_0\bm X_{lag}).
\]
Under the true partition,
\[
    \bm Y-\beta_0\bm X_{lag}
    =
    \bm Q_{K,\Lambda_0}\bm c_0+\bm r_0+\bm U,
\]
where \(\bm r_0\) stacks the regime-specific sieve approximation errors. Let \(\bm e_0=\bm M_{K,\Lambda_0}\bm r_0\).

\begin{lemma}[Known-partition replacement]
\label{lem:s-known-replacement}
Under Assumptions \ref*{ass:block-sieve} and \ref*{ass:oracle},
\begin{align}
    T^{-1/2}\sum_{t=1}^T\widehat Z_t(\beta_0,\Lambda_0)
    &=T^{-1/2}\sum_{t=1}^T U_tv_{t,0}+o_p(1),
    \label{eq:s-num}\\
    T^{-1}\sum_{t=1}^T\widehat Z_t(\beta_0,\Lambda_0)^2
    &=T^{-1}\sum_{t=1}^T U_t^2v_{t,0}^2+o_p(1),
    \label{eq:s-den}\\
    \max_{t\leq T}|\widehat Z_t(\beta_0,\Lambda_0)|&=o_p(\sqrt T).
    \label{eq:s-max}
\end{align}
\end{lemma}

\begin{proof}
Using Lemma \ref{lem:s-block-identities},
\[
    \widehat{\bm u}_0(\beta_0)
    =
    \bm M_{K,\Lambda_0}(\bm U+\bm r_0),
    \qquad
    \bm v_0=\bm M_{K,\Lambda_0}\bm w.
\]
Because \(\bm M_{K,\Lambda_0}\) is symmetric and idempotent,
\[
    T^{-1/2}\widehat{\bm u}_0(\beta_0)'\bm v_0
    =
    T^{-1/2}(\bm U+\bm r_0)'\bm v_0.
\]
By Assumption \ref*{ass:block-sieve},
\[
    T^{-1/2}\bm r_0'\bm v_0
    =
    T^{-1/2}\bm r_0'\bm M_{K,\Lambda_0}\bm w
    =o_p(1),
\]
which proves (\ref{eq:s-num}).

It remains to record the denominator and maximum bounds. Let \(G_T=T^{-1}\bm Q_{K,\Lambda_0}'\bm Q_{K,\Lambda_0}\). Since the rows of \(\bm Q_{K,\Lambda_0}\) have norm at most \(\zeta_K\) and the eigenvalues of \(G_T\) are bounded away from zero and infinity,
\[
    \|\bm v_0\|_\infty=O_p(1+\zeta_K).
\]
Moreover,
\[
    T^{-1}\bm Q_{K,\Lambda_0}'\bm r_0=O_p(a_K),
\]
and hence
\[
    \|\bm e_0\|_T\leq\|\bm r_0\|_T=O_p(a_K),
    \qquad
    \|\bm e_0\|_\infty=O_p\{a_K(1+\zeta_K)\}.
\]
The stated rates imply
\[
    \|\bm e_0\odot\bm v_0\|_T
    =O_p\{a_K(1+\zeta_K)\}=o_p(1),
    \qquad
    \|\bm e_0\odot\bm v_0\|_\infty
    =O_p\{a_K(1+\zeta_K)^2\}=o_p(\sqrt T).
\]
For the denominator, write \(\widetilde{\bm U}_0=\bm M_{K,\Lambda_0}\bm U\). Then
\[
    \widehat Z_t(\beta_0,\Lambda_0)
    =
    \widetilde U_{0,t}v_{t,0}+e_{0,t}v_{t,0}.
\]
Assumption \ref*{ass:oracle} gives
\[
    T^{-1}\sum_{t=1}^T
    \{(\widetilde U_{0,t}v_{t,0})^2-(U_tv_{t,0})^2\}=o_p(1),
\]
and Cauchy--Schwarz with \(\|\bm e_0\odot\bm v_0\|_T=o_p(1)\) gives the remaining cross and square terms. This proves (\ref{eq:s-den}). The maximum bound follows from
\[
    |\widehat Z_t(\beta_0,\Lambda_0)|
    \leq
    |U_tv_{t,0}|
    +
    |(\widetilde U_{0,t}-U_t)v_{t,0}|
    +
    |e_{0,t}v_{t,0}|,
\]
Assumption \ref*{ass:oracle}, and \(\|\bm e_0\odot\bm v_0\|_\infty=o_p(\sqrt T)\).
\end{proof}

\begin{lemma}[EL expansion]
\label{lem:s-el-expansion}
Let \(Z_{t,T}\) be scalar scores and \(\widehat V_T^Z=T^{-1}\sum_tZ_{t,T}^2\). Suppose zero lies in the interior of the scalar convex hull with probability tending to one and
\[
    T^{-1/2}\sum_t Z_{t,T}=O_p(1),
    \qquad
    \widehat V_T^Z\to_p\eta^2,
    \qquad
    \max_t|Z_{t,T}|=o_p(\sqrt T),
\]
where \(P(\eta^2\geq c)=1\) for some fixed \(c>0\). Then
\[
    2\sum_{t=1}^T\log(1+\widehat\lambda Z_{t,T})
    =
    \frac{\{T^{-1/2}\sum_tZ_{t,T}\}^2}
         {T^{-1}\sum_tZ_{t,T}^2}
    +o_p(1).
\]
\end{lemma}

\begin{proof}
The second-moment condition implies \(P(\widehat V_T^Z\geq c/2)\to1\). Let \(S_T^Z=\sum_tZ_{t,T}\) and \(Q_T^Z=\sum_tZ_{t,T}^2\). On this event, \(S_T^Z=O_p(\sqrt T)\) and \(Q_T^Z\geq cT/2\). Monotonicity of the scalar multiplier equation, evaluated at \(\pm 2|S_T^Z|/Q_T^Z\), gives
\[
    |\widehat\lambda|
    \leq
    2|S_T^Z|/Q_T^Z
    =O_p(T^{-1/2}).
\]
Since \(\max_t|\widehat\lambda Z_{t,T}|=o_p(1)\), a Taylor expansion of the multiplier equation gives
\[
    \widehat\lambda
    =
    \frac{\sum_t Z_{t,T}}{\sum_t Z_{t,T}^2}
    +o_p(T^{-1/2}).
\]
Substitution into the second-order expansion of the log empirical likelihood gives the displayed quadratic approximation.
\end{proof}

\begin{proof}[Proof of Theorem \ref*{thm:known-wilks}]
Apply Lemma \ref{lem:s-known-replacement} to reduce the feasible score to the matched block-projection oracle score \(U_tv_{t,0}\). Assumption \ref*{ass:oracle} gives the stable numerator limit, denominator consistency, feasible convex hull, and maximum-score condition. Lemma \ref{lem:s-el-expansion} then yields
\[
    \ell_T(\beta_0,\Lambda_0)
    =
    \frac{\{T^{-1/2}\sum_tU_tv_{t,0}\}^2}
         {T^{-1}\sum_tU_t^2v_{t,0}^2}
    +o_p(1)
    \Rightarrow\chi_1^2.
\]
\end{proof}

\section{Estimated-Partition Equivalence}

This section records the high-level likelihood-transfer argument behind Theorem \ref*{thm:estimated-equiv}. Let
\[
    S_T(\Lambda)=T^{-1/2}\sum_t\widehat Z_t(\beta_0,\Lambda),
    \qquad
    V_T(\Lambda)=T^{-1}\sum_t\widehat Z_t(\beta_0,\Lambda)^2.
\]

\begin{lemma}[Fixed-order estimated-partition likelihood transfer]
\label{lem:s-estimated-transfer}
Under the assumptions of Theorem \ref*{thm:known-wilks} and Assumption \ref*{ass:estimated-score},
\[
    \ell_T(\beta_0,\widetilde\Lambda)
    -
    \ell_T(\beta_0,\Lambda_0)
    =o_p(1).
\]
\end{lemma}

\begin{proof}
Lemma \ref{lem:s-known-replacement} and Assumption \ref*{ass:oracle} give \(S_T(\Lambda_0)=O_p(1)\), \(V_T(\Lambda_0)\to_p\eta_q^2\), and a true-partition maximum-score bound. Assumption \ref*{ass:estimated-score} gives
\[
    S_T(\widetilde\Lambda)-S_T(\Lambda_0)=o_p(1),
    \qquad
    V_T(\widetilde\Lambda)-V_T(\Lambda_0)=o_p(1),
\]
and the fixed-order estimated-partition maximum-score and convex-hull conditions. Lemma \ref{lem:s-el-expansion} applies to both arrays, and the two quadratic ratios differ by \(o_p(1)\) because both denominators are bounded away from zero with probability tending to one.
\end{proof}

\begin{proof}[Proof of Theorem \ref*{thm:estimated-equiv}]
By Lemma \ref{lem:s-estimated-transfer},
\[
    \ell_T(\beta_0,\widetilde\Lambda)
    =
    \ell_T(\beta_0,\Lambda_0)+o_p(1).
\]
The result follows from Theorem \ref*{thm:known-wilks}. If \(P\{\widehat q(\beta_0)=q_0\}\to1\), then \(\widehat\Lambda(\beta_0)=\widetilde\Lambda\) with probability tending to one. This proves Corollary \ref*{cor:qhat} under the fixed-order event version of Assumption \ref*{ass:estimated-score}.
\end{proof}
\section{Local Predictive Alternatives}

This section proves Theorem \ref*{thm:local-power}. Under the local alternative \(\beta_T=\beta_0+d_T\), the null-imposed residual satisfies
\[
    \bm Y-\beta_0\bm X_{lag}
    =
    \bm\mu+\bm U+d_T\bm X_{lag}.
\]
At the true partition, the block projection removes the regime-specific sieve approximation to \(\bm\mu\). The numerator replacement condition in Theorem \ref*{thm:local-power} gives
\[
    T^{-1/2}\sum_{t=1}^T \widehat Z_t(\beta_0,\Lambda_0)
    =
    T^{-1/2}\sum_{t=1}^T U_tv_{t,0}
    +
    T^{-1/2}\sum_{t=1}^Tg_{t,T}
    +o_p(1).
\]
By the local drift condition, the second term converges to \(\eta_q\delta_h\). The denominator replacement condition reduces the feasible denominator to the shifted oracle denominator. Moreover,
\[
    T^{-1}\sum_t(U_tv_{t,0}+g_{t,T})^2
    =
    V_T^{(q)}
    +
    2T^{-1}\sum_tU_tv_{t,0}g_{t,T}
    +
    T^{-1}\sum_tg_{t,T}^2.
\]
The last term is \(o_p(1)\) by the local square condition, and the cross term is \(o_p(1)\) by Cauchy--Schwarz and \(V_T^{(q)}=O_p(1)\). Thus the feasible denominator remains \(\eta_q^2+o_p(1)\). The maximum condition follows from the oracle maximum condition and the local maximum and replacement conditions. The feasible convex-hull condition and Lemma \ref{lem:s-el-expansion} give the known-partition limit \((Z+\delta_h)^2\). Lemma \ref{lem:s-estimated-transfer}, applied under the local-alternative estimated-score and convex-hull conditions, transfers the result to \(\widetilde\Lambda\). Correct order selection gives the unknown-order statement for \(\ell_T^{BA}(\beta_0)\).
\section{Implementation Details}

For simulations, each replication uses the same null-imposed profile path as the theory:
\[
    \beta
    \to
    \widehat q(\beta),\widehat\Lambda(\beta)
    \to
    \bm M_{K,\widehat\Lambda(\beta)}
    \to
    \widehat Z_t(\beta)
    \to
    \ell_T^{BA}(\beta).
\]
The proof object is global profile-sieve segmentation. Computation may use multiscale candidate generation, followed by final selection with the profile-sieve RSS and the information criterion in (\ref*{eq:qhat}). The implementation must report the candidate-generation method separately from the statistical estimator, because SBS, NOT, and related devices are computational accelerators rather than the main contribution.

For empirical work, the stable and break-adaptive statistics are reported together. The stable statistic documents conventional evidence of predictability; the adaptive statistic documents whether that evidence survives multiple nuisance breaks. The empirical section also reports \(\widehat q\), break dates, regime-specific nuisance fits, an estimated break-induced score component, and residualized predictor strength \(T^{-1}\bm w'\bm M_{K,\widehat\Lambda}\bm w\).

\end{document}
